\section{Problem Definition and Research Direction}\label{sec:problem-direction}

In a standard Digital Audio Workstation (DAW), the primary method of input is the ``Piano Roll''---a grid
where notes are
drawn manually across a timeline.
This approach is highly intuitive but fundamentally destructive from a structural perspective.
If a composer wishes to transpose a repeating 16-bar melody, change its rhythm, or conditionally alter its ending on the
final repetition, they must manually select and modify hundreds of individual note events.
The logical intent (e.g., ``repeat this phrase four times but play the snare on the off-beat during the last
repetition'') is lost, replaced by a static collection of data points.

To address this, programmers often turn to Algorithmic Composition.
However, existing domain-specific languages for music (such as Sonic Pi or Tidal) are largely tailored toward
\textit{live coding}.
They prioritize real-time performance, utilizing threaded execution models and continuous audio synthesis.
While powerful for improvisation, these environments introduce complexities related to thread
synchronization, execution latency, and audio engine configuration, which can be overwhelming for composers whose
primary goal is offline, deterministic arrangement.

Furthermore, general-purpose programming languages equipped with MIDI libraries (e.g., Python with \texttt{mido}) offer
offline generation but lack domain-specific abstractions.
The composer is burdened with calculating explicit delta-times and manually constructing low-level binary messages,
shifting the focus away from musical creativity and toward byte-level engineering.

\subsection{Problem Statement}\label{subsec:problem-statement}

There exists a gap in the current ecosystem of music programming tools: a lack of a domain-specific language optimized
for \textit{offline, structural composition} that produces standard, interoperable artifacts.
Specifically, the challenge lies in creating a system that allows composers to define musical ideas as reusable,
parameterized algorithms while keeping the resulting temporal structure deterministic, free of real-time
execution jitter,
and readable by many standard production tools.

\subsection{Proposed Solution}\label{subsec:proposed-solution}

This thesis proposes the design and implementation of a compiled domain-specific language
for musical
composition.
The system aims to bridge the gap between algorithmic logic and the rigid binary specifications of the MIDI protocol.

The proposed solution is organized around four design principles.

First, music is treated as a recursive hierarchy. The language abstracts musical motifs into functions (patterns) that
can be parameterized, reused, and nested, reflecting the true hierarchical nature of musical form.

Second, the source language uses a C-style imperative flow. By utilizing familiar control flow structures
(\texttt{for}, \texttt{loop},
\texttt{if}), Motivo allows composers to procedurally generate variations and repetitions without learning esoteric
functional paradigms.

Third, the compiler follows a zero-runtime compilation model. The compiler operates strictly offline.
It performs semantic validation and then uses a stateful lowering engine to completely unroll the composition's
logic into a flat timeline.
This keeps temporal resolution deterministic within the compiler pipeline without relying on an interpreter
during playback.

Fourth, MIDI is used as the semantic target. The compiler serializes the final timeline into a Standard MIDI File
(Format 1).
This supports practical interoperability, allowing the generated compositions to be imported into many DAWs
and MIDI-aware tools for
further orchestration and mixing.
